\section{Five Case Studies}
\label{sec:cases}

Each case study below is organized around three questions.
First, was the authority, on its statutory or constitutional face, intended as time-bounded?
Second, did the ratcheting pattern in fact occur,
in the sense that an authority granted as exceptional became routine or permanent?
Third, would the typed-rollback grammar developed in Part~\ref{sec:grammar} have structurally prevented the ratchet?
The third question is the most demanding and the one on which the Article's structural claim must stand or fall.
The Article is candid where the answer is contested or where the case study fits the framing less cleanly.
The five case studies vary in fit, and the variation is itself instructive.

\subsection{Weimar Article 48: the canonical case}
\label{sec:cases:weimar}

Article~48 of the Weimar Constitution \cite{ref_weimar_const_art48} authorized the Reich President
to take ``the necessary measures'' to restore public security and order
in cases where it had been ``substantially disturbed or endangered.''
The article required the Reich President to notify the Reichstag of measures taken,
and Article~48(3) permitted the Reichstag to demand the suspension of any such measure.
That suspension power, however, was politically conditional rather than a static institutional check:
it required a parliamentary majority willing to assemble against the sitting government, and across
1930 to 1933 the SPD's tactical \emph{Tolerierung} of Br\"uning withheld that majority, while
Hindenburg's dissolutions of the Reichstag in July 1930 and again in 1932 removed the occasion for a
suspension vote at precisely the moments it would otherwise have been most consequential.
The drafters of Article~48, working within the constitutional convention at Weimar in 1919,
appear to have understood the authority as a narrowly bounded emergency power
designed to provide the new republic with a tool for response to acute civil disorder
\cite{ref_caldwell_popular_sovereignty,ref_jacobson_weimar_pdf,ref_jacobson_schlink_weimar}.

The historical trajectory of Article~48's use is well-documented.
Through the 1920s, the authority was invoked in response to discrete crises:
the hyperinflation of 1923, the political assassinations of the early decade,
and the recurring municipal-level civil disturbances that characterized the era.
The invocations during the early period were generally followed by Reichstag review,
and the measures taken were, in many cases, subsequently regularized through ordinary legislation
or allowed to lapse \cite{ref_kershaw_hitler}.

The pattern shifted in the early 1930s, and the period 1930 to 1933 is best disaggregated rather than
treated as a single phase.
Heinrich Br\"uning's government (March 1930 to May 1932) pursued a legalistic, deflationary program of
emergency budgets enacted by Article~48 \emph{Notverordnung} after the Reichstag could not assemble a
majority for ordinary austerity legislation \cite{ref_bruening_notverordnung_1931}; the government
relied on Social Democratic \emph{Tolerierung} -- a tactical decision by the SPD parliamentary group not
to invoke Article~48(3) suspension -- and Hindenburg's dissolution of the Reichstag in July 1930
removed the immediate occasion for a suspension vote.
Franz von Papen's government (June to December 1932) took an openly anti-parliamentary turn, including
the July 1932 \emph{Preu{\ss}enschlag}, the Article~48-based deposition of the elected Prussian
government that became the subject of \emph{Preu{\ss}en contra Reich}.
Kurt von Schleicher's brief chancellorship (December 1932 to January 1933) attempted a \emph{Querfront}
coalition strategy that collapsed without producing comparable Article~48 measures.
Across the three governments Article~48's text -- which spoke of measures necessary to restore
public order -- had been stretched to cover budgetary decisions, currency policy, and the regulation of
commerce \cite{ref_caldwell_popular_sovereignty,ref_mommsen_weimar}.
The Reichstag Fire Decree of February 28, 1933 \cite{ref_reichstag_fire_decree_1933}, issued under
Article~48 in response to the burning of the Reichstag building, suspended civil liberties indefinitely
and authorized the \emph{Schutzhaft} that produced the immediate arrest of approximately eighty-one
Communist deputies and the harassment of Social Democratic deputies.
The Enabling Act of March 24, 1933 \cite{ref_enabling_act_1933} was procedurally distinct, but its
two-thirds threshold was reached because Article~48 had already been used to physically remove the KPD
opposition from the chamber; the path from Article~48 to the Enabling Act ran through the Reichstag Fire
Decree's preventive-detention provisions rather than through any generalized normalization of
emergency rule.

The historiography on whether the Article~48 ratchet was structural or political is contested
(the Kershaw/Mommsen/Caldwell historiographical contestation discussed in Part~\ref{sec:pattern}),
and the primary-source record is collected in Jacobson and Schlink's anthology
\cite{ref_jacobson_schlink_weimar} alongside Dyzenhaus's reconstruction of the
Kelsen-Schmitt-Heller debate \cite{ref_dyzenhaus_legality_legitimacy}.
The Article does not need to resolve this dispute.
It needs only the narrower observation that Article~48 contained no typed rollback obligation:
the Reich President could exercise the authority without exhibiting, at the moment of exercise,
a constructible path back to the prior state.
The decrees taken under the authority created facts on the ground -- suspended civil liberties,
imprisoned dissidents, dissolved political organizations -- that did not return to the prior state
when the formal emergency ended.

The constitutional order was not, however, wholly passive in the face of the trajectory.
\emph{Preu{\ss}en contra Reich}, decided by the Staatsgerichtshof on 25 October 1932
\cite{ref_preussen_contra_reich_1932}, was the case in which Article~48's limits were actually
litigated.
The Land of Prussia challenged the Papen government's July 1932 \emph{Preu{\ss}enschlag}, the
Article~48-based deposition of the elected SPD-Centre Prussian government and its replacement by a Reich
Commissioner.
The court partially upheld the Reich's authority to install a commissioner for limited functions but
partially struck down the deposition, holding that the Prussian ministers retained their constitutional
status as representatives of the Land in the Reichsrat and that the Reich could not, under Article~48,
displace them from that representative capacity.
The decision is direct evidence that the constitutional order resisted the ratchet in a way the
purely structural narrative elides: the limits of Article~48 were not only textually unspecified but
also judicially contestable, and the contest produced a partial substantive ruling against the executive
within months of the action.
The structural reading this Article advances accommodates that ruling: the absence of a typed
rollback witness made the ratchet structurally available without making it textually entailed, a
narrower claim than \emph{Preu{\ss}en contra Reich} disturbs.

The typed-rollback grammar would have structurally constrained Article~48's use in a way the original
text did not.
A grammar that required the Reich President, at the moment of exercising Article~48 authority,
to exhibit a witness for the path back to the prior state would have made decrees of that form
unconstructible, although the political pressure that produced the decrees would have remained.
The grammar would not have prevented the political use of Article~48 to circumvent the Reichstag,
but it would have prevented the conversion of that political use into permanent constitutional facts.
This is the cleanest application of the framing to a historical case,
and the case study where the structural absence is most legible.

\subsection{Section 230's incentive shadow and the DMCA 512 takedown regime}
\label{sec:cases:section230}

The framing of Section~230 of the Communications Decency Act \cite{ref_section230} as bearing
structural similarity to a delegated emergency authority is one this Article advances against the
consensus of the Section~230 academy, which has, in differing directions, resisted the characterization
of 230 as a takedown regime \cite{ref_goldman_section230_2019,ref_citron_wittes_bad_samaritans_2017,ref_keller_overremoval_2015}.
Eric Goldman has argued that 230 is best understood as a non-takedown immunity that decouples platform
incentives from content-moderation outcomes \cite{ref_goldman_section230_2019}; Danielle Citron and
Benjamin Wittes have argued for narrowing the shield on bad-Samaritan grounds rather than on
irreversibility grounds \cite{ref_citron_wittes_bad_samaritans_2017}; Daphne Keller's empirical work on
over-removal under intermediary-liability regimes argues against treating immunity withdrawal as a
takedown lever \cite{ref_keller_overremoval_2015}.
The Article's structural claim is therefore advanced with the academy's resistance in view, and the
characterization that follows separates the components of the regime more carefully than its prior
formulations did.

Section~230(c)(1) is a liability shield: it operates \emph{ex post}, when a plaintiff sues, and
immunizes a publisher from being treated as the speaker of content posted by another.
It does not authorize, require, or even mention takedown.
Section~230(c)(2) is structurally distinct from (c)(1).
The ``good Samaritan'' carve-out in (c)(2) explicitly immunizes good-faith action to restrict access to
material the provider considers obscene, lewd, lascivious, filthy, excessively violent, harassing, or
otherwise objectionable.
That sub-section is, on its face, a hedged authorization of takedown by the private intermediary, paired
with no statutory restoration obligation.
The Article treats (c)(2) honestly: it is already a takedown regime of a particular form, and the
absence of a structural rollback obligation within (c)(2) is one element of the pattern this Article
identifies.

The Digital Millennium Copyright Act's notice-and-takedown procedure
\cite{ref_dmca_512} is the operative notice-and-takedown regime in U.S. intermediary-liability law, and
unlike Section~230 it contains a structurally relevant partial rollback path: 17 U.S.C. \S~512(g)
obligates the service provider to restore content within ten to fourteen business days of a properly
formed counter-notice unless the complainant files suit \cite{ref_dmca_512g}.
The \S~512(g) counter-notice is, in the formal sense this Article develops, closer to a typed rollback
witness than either Section~230(c)(2) or the platform-self-regulatory takedowns that occur in the
broader shadow of immunity supplies; it is a procedurally required path back to the prior state, even if
the path is conditional on the user's willingness to incur the friction of filing.
The Article therefore does not claim that the intermediary-liability landscape lacks any
rollback-adjacent apparatus.
It claims, more narrowly, that Section~230's incentive shadow -- the (c)(1) immunity coupled with the
(c)(2) good-faith carve-out -- produces platform self-regulatory takedowns that operate outside the
DMCA's structured notice-and-counter-notice framework, and that those self-regulatory takedowns lack
the structural rollback witness that \S~512(g) supplies in the copyright-takedown adjacent regime.

The empirical record on takedown patterns is consistent with this narrower characterization.
Work on DMCA notice-and-takedown indicates that a substantial fraction of notices target content for
which the underlying infringement claim is dubious or absent \cite{ref_urban_quilter_dmca}, and that
the \S~512(g) counter-notice is rarely used because of the friction it imposes on the user whose
content was removed \cite{ref_seng_dmca}.
The counter-notice regime is, in other words, a typed rollback witness with high enough activation costs
that the underlying default approximates non-restoration; the platform's safe-harbor incentive favors
removal, and the counter-notice path is available but underutilized.
In the broader self-regulatory regime that operates in Section~230's incentive shadow, no analogous
witness exists at all: removals based on internal policy rather than statutory notice are not subject
to a \S~512(g)-style restoration path.

The typed-rollback grammar, applied to this regime, identifies the structural locus precisely.
The grammar does not displace \S~512(g) in the DMCA-counterpart context; rather, it generalizes
\S~512(g)'s logic and presses it back to the moment of removal.
A removal action under either the DMCA or the broader Section~230 self-regulatory shadow would, under
the grammar, require the platform to construct, at the moment of removal, a witness for restoration of
the content in the case that the underlying claim is later determined to be without merit.
The witness would have to be activated structurally rather than only on user objection.
A removal that cannot be reversed by construction would fail to type-check.

The application of the grammar to this regime is contested.
Goldman would observe that the grammar imposes a structural cost on intermediaries that Section~230
deliberately declines to impose; Keller would observe that the typed-rollback discipline does not
address the over-removal pattern her empirical work documents because it operates at the moment of
admission rather than at the moment of incentive formation; Citron and Wittes would argue, from a
different direction, that the discipline does not address the harms produced by content that should
have been removed but was not.
The Article acknowledges these objections without resolving them.
The structural claim is the narrower one: that Section~230's incentive shadow produces platform
self-regulatory takedowns that lack a structural rollback witness, even though DMCA \S~512(g) provides
a partial counter-example in the adjacent copyright-takedown regime.
The policy question of whether the grammar should apply to this regime is a separate question on which
the Article does not take a position.

\subsection{GDPR Article 17: erasure without rollback}
\label{sec:cases:gdpr}

Article~17 of the General Data Protection Regulation \cite{ref_gdpr_art17} grants data subjects a
right to erasure of their personal data in specified circumstances.
The article's enforcement has produced a body of erasure and de-referencing orders against search-engine
providers and platform operators, of which Case C-131/12, \emph{Google Spain v. AEPD},
ECLI:EU:C:2014:317 (13 May 2014) \cite{ref_google_spain} is the foundational example.
The erasure orders are time-bounded executive acts in the formal sense that they are issued by
administrative bodies operating under defined procedures, and they are subject to judicial review on
familiar administrative-law grounds.
The Court of Justice has, since \emph{Google Spain}, developed a substantial body of jurisprudence
operationalizing the balance the controller and the data-protection authority must strike.
Case C-136/17, \emph{GC and Others v. CNIL}, ECLI:EU:C:2019:773 (24 September 2019)
\cite{ref_gc_and_others_cnil_2019} clarified the scope of de-referencing for special-category data and
held that the balance between data-subject rights and public-interest access is re-evaluable in light of
changed circumstances.
Case C-507/17, \emph{Google LLC v. CNIL}, ECLI:EU:C:2019:772 (24 September 2019)
\cite{ref_google_cnil_2019} held that the right to de-referencing does not, as a matter of EU law,
extend to all domain extensions globally.
Case C-460/20, \emph{TU and RE v. Google LLC}, ECLI:EU:C:2022:962 (8 December 2022)
\cite{ref_tu_re_google_2022} addressed re-listing where allegedly inaccurate content is subsequently
shown to be accurate, treating the balance as one the controller must reassess.
The prior-regime precedent that established the threshold definition of personal-data processing that
the GDPR inherited is Case C-101/01, \emph{Lindqvist}, ECLI:EU:C:2003:596 (6 November 2003)
\cite{ref_lindqvist_2003}.

The ratcheting pattern in this regime takes a different form from the Weimar case.
There is no observed trajectory in which the erasure authority has been used to dissolve the underlying
constitutional order or to substitute for ordinary legislation.
The pattern that this Article identifies is the structural irreversibility of individual erasure orders
in the regulation's default operation.
An erasure order requires the controller to delete or de-index the targeted data.
The deletion is not, on the regulation's default operation, restored, although the underlying balance
is re-evaluable and the controller may re-index where the supporting facts change
\cite{ref_kuner_gdpr_erasure}.
The order is, in form, a time-bounded executive act.
It is, in effect, a permanent alteration of the data ecosystem absent affirmative reassessment by the
controller, who is under no structural obligation to re-evaluate at any defined interval.

The Article's framing here is narrower than in the Weimar case.
The claim is not that the GDPR Article~17 regime has ratcheted into something other than what it was
designed to be.
The claim is that individual erasure orders share the structural defect of an emergency authority
granted without a typed rollback witness: the order can be admitted without exhibiting, at the moment
of admission, a constructible path back to the prior state.
The CJEU's recognition that the underlying balance is re-evaluable in cases like \emph{GC and Others}
and \emph{TU and RE} establishes the doctrinal possibility of restoration in specific factual
configurations; it does not establish a structural witness for restoration at the moment of admission.
The witness is not currently a requirement of the regime.

The European Data Protection Board's Guidelines 5/2019 on the criteria of the Right to be Forgotten in
the search-engines cases \cite{ref_edpb_guidelines_5_2019} codify the controller's ongoing obligation
to revisit the balance, and they supersede the Article 29 Working Party Guidelines on Implementation of
the Court of Justice of the European Union judgment on ``Google Spain and Inc. v. Agencia Espa\~nola de
Protecci\'on de Datos and Mario Costeja Gonz\'alez'' (WP 225, 26 November 2014)
\cite{ref_a29_wp225_2014}.
The Article's claim that the regime is structurally irreversible at the moment of admission must be
read against these documents: the regime contemplates reassessment, but the reassessment is not a
typed rollback witness in the sense Part~\ref{sec:grammar} develops.

A more structurally proximate EU instrument is the Digital Services Act
\cite{ref_dsa_2022}, Regulation (EU) 2022/2065, whose principal obligations applied to hosting
providers from 17 February 2024 and to designated Very Large Online Platforms and Search Engines
from 25 August 2023.
Article~16 installs a notice-and-action regime requiring providers to receive, to process in a timely
and non-arbitrary manner, and to decide notices of allegedly illegal content; Article~17 attaches a
statement-of-reasons obligation to any restriction imposed on recipients; Article~20 requires an
internal complaint-handling system that must remain available for at least six months after the
contested decision and that must operate electronically and free of charge; and Article~21 supplies
certified out-of-court dispute settlement through bodies recognized by national Digital Services
Coordinators.
The conjunction of these provisions, and the six-month Article~20 complaint window in particular,
comes structurally closer to what this Article calls a typed rollback witness than any regime examined
above: a removal cannot, in the DSA's regime, be admitted without an accompanying procedural path by
which it may be reversed.
The DSA's existence strengthens the structural case advanced here.
The European Union has demonstrated that the construction is buildable and operative in a regulated
market, and the regime's confinement to intermediary content moderation illustrates the
generalization gap this Article argues should be closed.

The data-protection academic literature has framed the right structurally rather than
transactionally, but none of the leading treatments proposes a typed-rollback discipline as the
corrective.\footnote{Orla Lynskey's \emph{The Foundations of EU Data Protection Law} (Oxford 2015)
\cite{ref_lynskey_foundations_eu_data_protection_2015}, particularly Chapters 6 and 7, treats the
right to be forgotten as a structural element of the data-protection regime rather than as a remedy
attached to particular harms; Lee A. Bygrave's \emph{Data Privacy Law: An International Perspective}
(Oxford, 2d ed. 2014) \cite{ref_bygrave_data_privacy_law_2014} situates the EU architecture in its
comparative-law context; Mantelero's writing on GDPR enforcement
\cite{ref_mantelero_gdpr_enforcement} and Kuner's analysis of the right to be forgotten
\cite{ref_kuner_gdpr_erasure} both observe the structural irreversibility of individual erasure
orders in the regulation's default operation.}
The Article's contribution to this literature is the identification of the structural similarity to the
emergency-authority pattern, and the proposal that the typed-rollback grammar provides a more
principled corrective than ad hoc restoration mechanisms.

The fit of the case study to the framing is intermediate.
The GDPR Article~17 regime is not on its face an emergency-authority regime, and the data-protection
scholars who have written on it have not, in general, framed it as such.
The structural similarity is at the level of the form of authority rather than the substance of the
regime.
The Article includes the case study because the structural similarity is real and because the
corrective the grammar proposes is applicable to the regime even though the regime is not, in its
self-understanding, an emergency-authority regime.

\subsection{FISA Section 702: programmatic acquisition and its emergency sub-provision}
\label{sec:cases:fisa702}

Section~702 of the Foreign Intelligence Surveillance Act, codified at 50 U.S.C. \S~1881a
\cite{ref_fisa_702_statute}, is a programmatic certification regime for the acquisition of
foreign-intelligence information from non-U.S. persons reasonably believed to be located outside the
United States; the program is reviewed annually by the Foreign Intelligence Surveillance Court under
50 U.S.C. \S~1881a(j).
Section~702 is not, in itself, an emergency authority.
The seventy-two-hour mechanism in 50 U.S.C. \S~1881a(c)(2) is one sub-provision, permitting the Attorney
General and Director of National Intelligence to determine targeting prior to FISC approval of a new
certification.
The case study this Article advances is about that sub-provision's interaction with the broader
\S~702 retention and querying framework, as documented in declassified FISC opinions and PCLOB reports,
rather than about Section~702 as a whole being an emergency authority.
Reading the case study otherwise would conflate the programmatic certification regime with the narrower
within-section emergency mechanism, and the Article disclaims that conflation at the outset.

The operative public record on the program runs through a sequence of declassified FISC opinions and
PCLOB reports.
\emph{In re DNI/AG 702(g) Certifications}, FISC Mem. Op. (Bates, J., Oct. 3, 2011), declassified August
2013 \cite{ref_fisc_bates_2011_upstream}, was the upstream-collection multiple-communications-transaction
opinion that found a Fourth Amendment violation in NSA's then-current upstream practices and required
minimization changes.
\emph{Memorandum Opinion and Order}, FISC (Collyer, J., April 26, 2017)
\cite{ref_fisc_collyer_2017_about} disclosed compliance failures in upstream ``about'' collection and
led NSA to halt about-collection.
\emph{Memorandum Opinion and Order}, FISC (Boasberg, J., Oct. 18, 2018), declassified 2019
\cite{ref_fisc_boasberg_2018} addressed querying-procedure compliance.
\emph{Memorandum Opinion and Order}, FISC (Contreras, J., April 21, 2022), declassified May 2023
\cite{ref_fisc_contreras_2022} addressed FBI query compliance and was extensively cited in the
2023--24 reauthorization debate.
The Privacy and Civil Liberties Oversight Board's three reports across the period are the PCLOB
\emph{Report on the Surveillance Program Operated Pursuant to Section 702 of the Foreign Intelligence
Surveillance Act} (July 2, 2014) \cite{ref_pclob_702_2014}, the \emph{Recommendations Assessment
Report} (January 2017) \cite{ref_pclob_recommendations_2017}, and the \emph{Report on the Section 702
Surveillance Program} (September 2023) \cite{ref_pclob_702_2023}.
The reauthorization legislation that re-entered the structural picture in this period comprises the
FISA Amendments Reauthorization Act of 2017, Pub. L. No. 115-118 (signed January 2018)
\cite{ref_faa_reauth_2017}, and the Reforming Intelligence and Securing America Act, Pub. L. No.
118-49 (April 2024) \cite{ref_risaa_2024}, the latter of which tightened FBI querying procedures and
audit requirements.

The structural defect, on the framing this Article develops, is the absence of a typed rollback witness
for information collection at the admission step.
The seventy-two-hour mechanism authorizes the collection;
it does not require, at the moment of authorization, the construction of a path back to the
pre-collection state.
The targets of the collection do not return to a pre-collection state when the seventy-two-hour
authorization is later reviewed or when the underlying certification is modified.
The collected information persists in agency holdings.

Existing law contains rollback-adjacent mechanisms: FISC-ordered destruction of unlawfully acquired
material under 50 U.S.C. \S~1881a(i)(3), the 50 U.S.C. \S~1809(a) criminal bar on use or disclosure of
unlawfully acquired information, the five-year retention defaults in approved minimization procedures,
and the post-2018 querying procedures.
These mechanisms are partial: they operate at the audit and lapse steps rather than at the admission
step, they do not bind the upstream collection, and they are not constructively reversible witnesses in
the sense Part~\ref{sec:grammar} develops.
The structural claim this Article advances is not that the regime contains no reversibility apparatus,
but that the apparatus it contains is not of the form a typed rollback witness would supply.

The case-law backdrop frames the litigation posture of the program.
\emph{Clapper v. Amnesty International USA}, 568 U.S. 398 (2013) \cite{ref_clapper_v_amnesty_2013},
established the standing doctrine that has structured almost all \S~702 litigation in Article~III
courts; the speculative-injury framework Clapper announced makes the structural-irreversibility claim
non-justiciable in the form in which the Article's framing might initially suggest it should be
litigated.
\emph{United States v. Hasbajrami}, 945 F.3d 641 (2d Cir. 2019) \cite{ref_us_v_hasbajrami_2019}, is
the principal federal-appellate merits decision on \S~702 querying of U.S.-person information; the
court treated the querying of incidentally collected U.S.-person communications as a distinct Fourth
Amendment event from the upstream acquisition, a posture directly relevant to whether the typed-rollback
grammar would apply at the querying step, at the acquisition step, or at both.

The Article's claim here is the most cautious of the five case studies.
The public record on Section~702 acquisition and the seventy-two-hour sub-provision is incomplete.
The declassified opinions and PCLOB reports provide an authoritative but partial picture of the
program's operation.
A national-security-law scholar with classified-clearance access to the full record may have grounds to
dispute the characterization above.
The structural claim is offered as a working hypothesis on the basis of the available public record,
with explicit acknowledgement that the record is incomplete.
The typed-rollback grammar's applicability to surveillance authorities depends on whether information
collection can, in principle, be made reversible by construction.
Whether such a construction is operationally feasible in the surveillance context, including in light
of the cryptographic literature on revocable encryption \cite{ref_revocable_encryption_lit}, is an
empirical question on which the Article does not take a position.

\subsection{The 2001 AUMF: a quarter-century of unbounded authorization}
\label{sec:cases:aumf}

The Authorization for Use of Military Force enacted by Congress on September 18, 2001 \cite{ref_aumf_2001}
permits the President to use ``all necessary and appropriate force'' against ``those nations, organizations, or persons
he determines planned, authorized, committed, or aided'' the September 11 attacks.
The text contains no geographic limit, no sunset, and no requirement of periodic reauthorization.

The trajectory of the AUMF's use is documented at length by the Brennan Center
\cite{ref_brennan_center_aumf} and by scholarly accounts in Chesney \cite{ref_chesney_aumf}.
By 2024, the authorization had been cited as legal basis for armed action in at least eight named
theaters, including territories of states not in existence as targets of policy at the time of the
authorization's enactment.
The running list is documented in the Department of Defense's 2018 Section~1264 Report on the legal and
policy frameworks guiding the use of military force \cite{ref_dod_1264_report_2018}, which named seven
countries (Afghanistan, Iraq, Syria, Yemen, Somalia, Libya, and the Philippines) and identified Niger
and Kenya in earlier transmittals; in the Congressional Research Service's continuing analysis of the
2001 AUMF and its application \cite{ref_crs_r43983_aumf}; and in the periodic War Powers Resolution
reports transmitted to Congress under 50 U.S.C. \S~1543
\cite{ref_war_powers_reports_2021_2023}.
Jameel Jaffer's compilation of the public legal record on drone operations
\cite{ref_jaffer_aumf} supplies the legal-theory critique of the associated-forces doctrine through
which the authority has been extended.
The Islamic State of Iraq and the Levant, which did not exist in any form in 2001, has been
characterized as a covered organization under the AUMF on the legal theory that it is a successor
entity to al-Qaeda in Iraq.
The expansion of the authority through the doctrine of associated forces has been documented in
Department of Justice white papers, declassified Office of Legal Counsel opinions, and the periodic war
powers reports.

The pattern admits both a political and a structural reading, and the structural reading must be
advanced carefully because the political record is not one of legislative inertia.
A run of bipartisan repeal-and-replace proposals demonstrates that the political coalition for repeal
has at times assembled.
The Kaine-Young AUMF repeal-and-replace framework was introduced as S.1228 in the 116th Congress and
reintroduced through successive sessions \cite{ref_kaine_young_aumf}; the Lee-Murphy proposals
addressed parallel concerns; the Corker-Kaine 2018 framework attempted a successor-authority structure
with sunset and reporting requirements \cite{ref_corker_kaine_2018}; and in 2023 Congress in fact
repealed both the 1991 Authorization for Use of Military Force Against Iraq and the 2002 Authorization
for Use of Military Force Against Iraq, which President Biden signed, although the 2001 AUMF itself was
left intact \cite{ref_iraq_aumf_repeal_2023}.
What has not converged is agreement on a successor authority for the 2001 AUMF specifically.
The structural reading this Article advances operates alongside the political account rather than
displacing it: the original grant carries an additional feature that lowers the cost of continuation
relative to replacement at each successive period, and the typed-rollback grammar addresses that
asymmetry.

The typed-rollback grammar's application to the AUMF is the most politically charged of the five case studies.
The application is structural: the original AUMF was enacted without a typed rollback witness,
which is to say, without a requirement that the President exhibit, at the moment of authorizing the use of force,
a constructible path back to the pre-authorization state.
The grammar would have required the original authorization to specify the conditions under which the use of force
was to terminate and the state of affairs to which the post-termination order would return.
A military action that could not exhibit, at the moment of authorization, a constructible path back to the prior state
would have failed to type-check.

The political content of this case study is not the Article's subject.
The Article does not take a position on the merits of any particular use of force under the AUMF or on
the desirability of any particular successor authorization, and the structural reading is supplementary
to, not a substitute for, the political account that Kaine-Young, Corker-Kaine, and the 2023 Iraq AUMF
repeal record represent.
The structural claim is the narrower one: that the original AUMF's grant of unbounded authority is the
instance in which the structural absence is most visible, and that the typed-rollback grammar would
have made the ratchet structurally less available by raising the cost of continuation relative to
replacement.
A reader who disagrees with the structural claim on this case study can disagree with it without
disturbing the Article's claims about the other four cases.
The AUMF is included because it is the most prominent contemporary American example of the ratcheting
pattern and because a structural account that omitted it would be incomplete; it is also the case study
where the political account is strongest, and the Article concedes that the structural reading must
ride alongside the political one rather than displace it.

\subsection{Comparative summary}

Table~\ref{tab:cases} summarizes the five regimes against the four-component grammar of Part~\ref{sec:grammar}. The column \emph{Rollback Witness} reports whether a typed rollback witness is structurally available at the moment of admission; it is the structural feature this Article identifies as the missing element across all five regimes.

\begin{table}[h]
\centering
\caption{Comparative summary of the five case studies. The structural pattern is the absence of a typed rollback witness at the moment of admission; the ratchet outcome varies with political and institutional context but is structurally available in each regime.}
\label{tab:cases}
\small
\begin{tabular}{@{}p{0.13\linewidth}p{0.18\linewidth}p{0.20\linewidth}p{0.15\linewidth}p{0.22\linewidth}@{}}
\toprule
Regime & Authority Form & Return Mechanism & Rollback Witness & Ratchet Outcome \\
\midrule
Weimar Art.~48 & Emergency decree & Reichstag suspension under Art.~48(3) & Absent & Absorbed into Nazi state \\
\addlinespace
2001 AUMF & Military force authorization & Repeal-and-replace (failed in successor convergence) & Absent & Twenty-five years, eight named theaters \\
\addlinespace
FISA \S~702(c)(2) & Seventy-two-hour emergency targeting sub-provision & FISC review, minimization procedures, querying audits & Partial: purge orders under \S~1881a(i)(3); retention defaults & Collection persists; querying audit added post-2018 \\
\addlinespace
GDPR Art.~17 & Erasure order against intermediary or controller & CJEU re-evaluability; controller re-indexing on changed facts & Partial: de-referencing reversible in principle, not by default & Default-restoration absent; re-evaluation rare in practice \\
\addlinespace
\S~230 \&\ DMCA \S~512 & Liability shield (230(c)(1)/(c)(2)) plus notice-and-takedown (512) & DMCA \S~512(g) counter-notice within ten to fourteen days & Partial: counter-notice in copyright; absent in Section 230 incentive shadow & Platform self-regulation entrenched as default takedown regime \\
\bottomrule
\end{tabular}
\end{table}
